\scanpage{235}
\uline{\mspage{2}.}
Plus généralement, sur un espace \unsure{top} \inserted{annelé} $X$, soient $K',K''$ deux complexes de Modules, et
$\mathcal G$ un faisceau abélien

$\mathcal P$ un $\mathcal G$-torseur à droite
\diagram{\begin{tikzcd}
\mathcal P\arrow[r,"u",no head]&\unsure{Z^0}\underline{\operatorname{Hom}}^\bullet(K',K'')\\
\mathcal G\arrow[r,"v"]&\underline{\operatorname{Hom}}^\bullet(K'',K'')\\
\mathcal G\arrow[r,"\unsure{h|k}"]&\underline{\operatorname{Hom}}^{-1}(K^{\unsure{\prime|\prime\prime}},K'')
\end{tikzcd}}
\[
\begin{aligned}
u(pg)&=u(p)v(g)\\
v(g+g')&=v(g)+v(g')\\
u(g)&=d\unsure{h|k}(g)+\unsure{h|k}(g)d
\end{aligned}
\]
On veut en conclure un hom dans la cat. dérivée \inserted{$K'\to K''$} d'où
\[
R\Gamma_X(K')\longrightarrow R\Gamma_X(K'')
\]
(compatible avec la structure sur $\deleted{\unsure{R\Gamma_X}}\Gamma_X(\underline O_X)$ ....)
[et plus gén. en relativisant par rapport à un morphisme d'espaces annelés $X\unsure{\mathbin{-}|\to}T$ ...]

On introduit l'extension $\overline{\mathcal G}$ de $\mathbf Z$ par $\mathcal G$
\diagram{\begin{tikzcd}
0\arrow[r,no head]&\mathcal G\arrow[r,no head]&\overline{\mathcal G}\arrow[r,"j",no head]&\mathbf Z\arrow[r,no head]&0
\end{tikzcd}}
définie au moyen de $\mathcal P$ par la relation
\[
j^{-1}(1)\simeq\mathcal P\qquad(\text{isom de }\uline{$\mathcal G$-torseurs}\text{ ...})
\]
Soit
\[
\mathcal H=\underline{\operatorname{Hom}}^\bullet(K',K'')
\]
et soit $L$ le complexe
\diagram{\begin{tikzcd}
0\arrow[r]&0\arrow[r]&\overset{\text{``}L^0}{\mathcal G}\arrow[r,no head]&\overset{\text{``}L^1}{\overline{\mathcal G}}\arrow[r]&0\arrow[r]&0
\end{tikzcd}}
\margin{234}
